\section{Reference Implementation}
\label{sec:implementation}

\subsection{Package layout}

The MIT-licensed reference implementation targets Python 3.10 or later and
PyTorch 2.10 or later. \Cref{tab:module-map} maps the specification to the
released modules.

\begin{longtable}{@{}p{0.32\textwidth}p{0.61\textwidth}@{}}
\caption{Implementation map.}
\label{tab:module-map}\\
\toprule
Module & Responsibility \\
\midrule
\endfirsthead
\toprule
Module & Responsibility \\
\midrule
\endhead
\texttt{config.py} & Validated immutable configuration objects and branch
schedules. \\
\texttt{state.py} & External state dataclasses, detach utilities, and logical
tensor-payload accounting. \\
\texttt{gated\_delta.py} & Causal stem, Gated Delta Rule-2 token recurrence,
write-surprise signal, sequence and step paths. \\
\texttt{local\_attention.py} & Rolling exact local KV state and causal
softmax read. \\
\texttt{exact\_cache.py} & Deterministic fixed-capacity priority insertion
and exact projected-KV attention. \\
\texttt{episodic.py} & Current/live/archive life cycle, slot aggregation,
router, and archive ring. \\
\texttt{compression.py} & Thin-QR small-core SVD factor truncation and
discarded-tail accounting. \\
\texttt{sketch.py} & Reproducible signed CountSketch maps and TensorSketch
features. \\
\texttt{model.py} & Block scheduling, branch fusion, residual path, language
model head, loss, and recurrent API. \\
\texttt{mqar.py} & Small ALUCLU diagnostic MQAR generator and metric. \\
\texttt{zoology\_mqar.py} & Historically aligned Zoology MQAR semantics and
masked loss. \\
\texttt{protocols/*.json} & Immutable protocol manifests, including the
historical mismatch record and a separately identified repair. \\
\bottomrule
\end{longtable}

\subsection{Configuration validation}

External dimensions and capacity budgets are validated before tensor
allocation. The configuration rejects non-positive dimensions, incompatible
head splits, invalid feature bounds, even sketch-replica counts where a median
requires an odd count, and inconsistent scheduling values. These checks are
not cosmetic: a silently truncated head dimension or non-positive feature
floor would invalidate the equations in \Cref{sec:memory-lanes}.

The implementation intentionally does not force \(r+S<d\). Configurations
with \(r+S\ge d\) remain mathematically valid, but the small SVD core can grow
to \(d\times d\), losing the intended compression-time advantage. This is
reported as a performance condition rather than prohibited as a semantic
error.

\subsection{Numerical precision}

The recurrent fast-weight matrix, token-dependent decay, and Gated Delta
update execute in fp32 even when the surrounding model uses fp16 or bfloat16.
Local and exact-cache softmax probabilities are computed in fp32 before being
cast back to the query dtype. QR and SVD promote fp16 or bfloat16 inputs to
fp32, then cast retained factors and the reported discarded tail back to the
input dtype.

These promotions protect the most sensitive reductions, but they do not make
the whole architecture numerically exact. Episodic persistent factors,
\(z\), route sums, sketches, and the error ledger remain in model dtype.
Furthermore, the ledger covers algebraic rank-truncation tails; it does not
add an explicit term for floating-point QR/SVD error or down-cast rounding.

\subsection{State ownership and graph control}

State dataclasses are treated as values at the public boundary. Each call
returns a new logical state rather than mutating a hidden global cache. This
supports interleaved streams and makes state transfer visible to callers. It
also means the caller owns lifecycle decisions:

\begin{itemize}
  \item pass \texttt{None} to begin a new stream;
  \item carry returned state to continue a stream;
  \item call \texttt{model.detach\_state(state)} at a training boundary when
  the old autograd graph must be released;
  \item do not reuse one sample's state for an unrelated sample.
\end{itemize}

The v0.1.0 API does not include a per-example reset mask for ragged batches and
does not define a versioned state serializer. Production serving would need
both before arbitrary stream regrouping or durable cache persistence.

\subsection{Determinism and discrete cache selection}

Exact-cache priority selection is batch-local and deterministic for a fixed
input. Empty slots are filled in index order. Once full, the first minimum
selected by the tensor reduction is replaced only if the incoming priority is
strictly greater. Equal priorities preserve the older record.

The priority is detached before discrete selection. Therefore cache-read loss
does not differentiate through the ranking decision into the Gated Delta
write residual. Accepted key/value projections can receive gradient through
the read; rejected records do not receive a key/value write contribution from
that step. This boundary is an intentional implementation fact, not a claim
that hard selection is optimal.

\subsection{Protocol provenance}

MQAR is particularly sensitive to target placement and next-token shifting.
The repository therefore assigns distinct protocol identifiers to:

\begin{enumerate}
  \item the small ALUCLU next-token diagnostic generator;
  \item the target semantics reconstructed from the historical Zoology
  generator;
  \item the historical raw manifest with its documented layer/config-count
  mismatch;
  \item the smallest runnable repair, under a new identity.
\end{enumerate}

This prevents a repaired configuration from silently inheriting the identity
of a non-runnable historical command and prevents two different target
alignments from being reported as the same benchmark.

\subsection{Testing strategy}
\label{sec:test-strategy}

The 50-test suite emphasizes invariants that can detect semantic drift:

\begin{itemize}
  \item dense equivalence of factor compression at sufficient rank;
  \item agreement between discarded singular tail and best-rank Frobenius
  error in float64 tests;
  \item error-ledger accumulation across repeated merges;
  \item exact \(z\) denominator preservation after factor compression;
  \item router identity at zero temperature and denominator-mass conservation;
  \item local, episodic, and exact-cache causality;
  \item full scan, multiple chunkings, and token-step equivalence;
  \item strict top-\(C_x\) cache retention and fixed initialized cache shape;
  \item state saturation after sufficient streaming tokens;
  \item archive capacity, compression depth, and ring eviction behavior;
  \item finite outputs, losses, and gradients on small random examples.
\end{itemize}

Tests establish implementation invariants, not model intelligence. In
particular, a passing cache-retention test proves that the policy keeps the
largest priorities, not that priority predicts future utility.

\subsection{Reference-backend performance characteristics}

The current backend uses a Python loop over tokens. Exact-cache updates employ
full-shape \texttt{torch.where}; episodic transitions create tuples and
concatenate factors; archive boundaries invoke library QR/SVD routines. These
choices favor clarity and direct comparison with the specification. They
should be treated as a correctness oracle and CPU fallback.

A production accelerator path should preserve the same externally tested
semantics while introducing:

\begin{enumerate}
  \item chunk-parallel or scan kernels for Gated Delta recurrence;
  \item fused causal stems, projections, recurrence, and output gating;
  \item ring-buffer local KV storage without repeated concatenation;
  \item in-place indexed cache writes with a compact selection reduction;
  \item packed episodic tensors with explicit masks instead of Python tuples;
  \item asynchronous or scheduled capsule compression;
  \item profiler-guided fusion constrained by numerical parity tests.
\end{enumerate}

Kernelization is a separate engineering phase because a fast kernel that
changes replacement ties, write-then-read order, fp32 accumulation, or
segment-boundary transitions would implement a different model.

\subsection{Minimal recurrent use}

\begin{lstlisting}[language=Python,caption={One-token recurrent inference.},
label={lst:minimal-step}]
import torch
from aluclu import AlucluLanguageModel, build_research_config

config = build_research_config(
    d_model=64,
    n_layers=4,
    local_window=64,
    segment_size=16,
    live_segments=16,
    archive_slots=4,
    archive_segments_per_slot=4,
    archive_rank=8,
    exact_cache_capacity=64,
)
model = AlucluLanguageModel(vocab_size=8192, config=config).eval()

state = None
token = torch.tensor([7])
logits, state = model.step(token, state)
\end{lstlisting}

The explicit \texttt{state} variable in \Cref{lst:minimal-step} is essential:
dropping it converts every call into a one-token independent sequence.
